\subsection*{数据库提交确认与消费确认分别丢失}

Q 已向数据库发出 COMMIT，连接中断，消费者不知道是否提交。恢复者连接原履约库，按 I1 查询：完整终态、F1 与回报都存在，就恢复原结果。都未提交则沿原 I/B 重试。受控入口下若只剩其中一半，转完整性核查。这里未知的是数据库已经作出的提交决定。

另一条历史是 Q 的提交确认明确到达，消费者再向 RabbitMQ 发送 ACK，ACK 丢失或进程退出。代理可能保留未确认投递。channel 或 connection 关闭并被检测到以后，手工确认模式下的未确认投递会重新入队，检测和恢复需要时间~\cite{rabbitmq-confirms}。新消费者收到后按 I1 恢复 F1，并用这次投递所属 channel 和 tag 确认。

\CTenFigure{ch10-original06}{数据库提交确认丢失与代理 ACK 丢失发生在两个独立位置}{fig:ch10-two-unknowns}

消费者把 tag 与本次投递的 channel 一起保存，在这个 channel 上按有效未确认集合发送确认。重复、未知或跨 channel 的确认会引发协议异常并关闭 channel~\cite{rabbitmq-confirms}。连接恢复后消息是否还可交付，还依第九章已经写明的队列存续和保存条件。

这些条件继续保持：消息通过默认交换机路由到指定目标 \code{fulfillment.paid}。拓扑版本为 \code{fulfillment-v1}，使用持久消息与持续存在的持久 quorum 队列。关闭消息/队列 TTL、自动删除与交付次数自动丢弃，容量压力拒收新发布。quorum 的默认交付上限和默认 \code{at-most-once} 死信转交可能使副本离开恢复路径~\cite{rabbitmq-quorum-queues}。若实际部署采用有限保留，消费方保存可关联的完成证据，上游按原 EV 核对未完成项，并在授权规则下补发原输入。发布 \code{accepted} 记录代理当时的受理事实，后续消费责任通过这套完成证据与补发记录继续衔接。

业务暂时不可处理时，保留未确认责任并限制在途量，等待数据库恢复。需要长时间等待、兼容读取或人工决定的输入，把完整来件及后续责任持久转入核查，再确认原投递。后续工作有负责人、扫描位置、重试时间和授权决定。持续立即 nack 后重入队，可能形成反复投递的 CPU 与网络循环~\cite{rabbitmq-confirms}。调度通过有界并发与延后运行，为数据库恢复和核查处理留出机会。

\subsection*{RabbitMQ 批量确认也会跨过空洞}

本章主方案明确调用 \code{basic.ack(delivery\_tag, multiple=false)}，只确认指定投递。设同一 channel 的 tag 7 尚未提交，tag 8 已经完成 Q。确认 8 时使用 \code{false}，7 继续保留在未确认集合，稍后可以完成或恢复。

若为了减少网络帧改用 \code{multiple=true}，代理会确认该 channel 截至指定 tag 的全部未确认投递~\cite{rabbitmq-confirms}。ACK 8 会连同 7 一起确认。进程随后退出，7 的业务尚未提交，重投责任却已经结束。批量确认需要维护当前 channel 的未确认集合和连续完成边界，只越过已经具备持久处理依据的投递。另一 channel 的 tag 独立计算。

这里的完成依据也可以是持久隔离交接。完整输入、原因和继续责任已经保存，代理责任就可以移交，消费方自己的待办保持开放。处理结果标明此次转入核查，业务履约则在取得相应依据以后另行推进。

\subsection*{履约回报与实际发货怎样继续}

先把 RabbitMQ 主线继续到购买库：Q 中的回报待办使另一发布器持续发送 \code{fulfillment.\allowbreak accepted}。购买服务收到回报，按报告身份、O1 与 F1 核对原归属。随后它在购买库自己的事务中保存接收记录，并共同更新 \code{orders} 表中的履约进度字段 \code{fulfillment\_state}。回报确认丢失，重复接收恢复同一迁移。此时两库最终可以达成相应进度。购买库状态变化发生在它自己的消费事务中，履约库 Q 先保存回报待办，为这次迁移提供持续输入。

第 7.4 节的订单结构只允许 \code{unfulfilled}。启用本节回报入口以前，先显式扩展检查约束，增加“履约库已受理”的阶段。受控入口验证回报后，才把该订单推进为 \code{fulfillment\_accepted}，重复回报读取原接收结果。这个阶段证明 F 已建立，发货事实由后面的独立步骤确定。

\begin{BookCode}
ALTER TABLE orders DROP CONSTRAINT orders_fulfillment_state_check;
ALTER TABLE orders ADD CONSTRAINT orders_fulfillment_state_check
  CHECK (fulfillment_state IN ('unfulfilled','fulfillment_accepted'));
\end{BookCode}

实际发货再跨到仓储或承运商。如果先保存 I.done，再去调用发货，调用以前退出就会留下建单成功、尚未发货。如果先发货，返回结果丢失以后重试，也可能重复外部动作。F 的存在为调度提供持久起点，发货操作保存自己的 dispatch 身份、原输入、固定下游身份、尝试和结果证据。

下游支持稳定操作身份及结果查询时，调度者按原 dispatch 身份恢复，沿第七、八章的规则处理未知结局与接管。消费重投或租约换代以后，调度者从原映射取得同一发货键。下游只有可能重复生效的调用、又没有足够排他和查询证据时，未知结果进入持久核查。初次建单的不变量在 Q 内成立，外部包裹次数则依据仓储或承运商实际执行的规则判断。


\subsection*{横向比较：Kafka 的分区位置怎样推进}

前面已经从 RabbitMQ 投递走到 Q，再接上购买库回报和独立发货。下面保持相同 I/B 业务仲裁，只把输入改为第九章介绍的 Kafka 日志，比较其组进度、批量完成边界和事务输出。这里的组位置用于这套替代输入模型，RabbitMQ 主方案继续使用 channel 与投递确认。

换到第九章的 Kafka 日志，先取主题 \code{order-events} 的分区 2，其中已有 offset 100、101、102 三条记录。offset 是记录在该分区中的位置。履约订阅使用 Kafka 4.1 的传统 consumer group，记为 G：多个消费者实例以同一组身份共同承担一份订阅，各分区在组内分配给一个实例读取。设 G 有 W1、W2 两个实例，当前分区 2 分配给 W1，W2 尚未取得该分区。G 为分区 2 持久保存的下一读取位置记为 \code{next=100}~\cite{kafka-design,kafka-consumer-api}。4.1 另有预览的 share group，本节推导限定传统组。

W1 用 \code{poll} 从分区取回记录，当前内存位置随读取推进；持久位置只有经过提交才改变。提交 offset 表示恢复时下一条待处理记录的位置，例如提交 101 表示从位置 101 继续。W1 退出后，分区可以重新分配给 W2，W2 在原组进度保留、所需日志可读的条件下，从 G 的已提交位置恢复~\cite{kafka-consumer-api}。

消费位置与日志保留各自管理。G 提交 103，只推进 G 的恢复起点，日志中的 100 至 102 按主题保留策略继续存放或清理。若独立审计组订阅同一主题，它拥有自己的分区位置，可以读取同一条记录。这里消费方提交的是一份订阅的分区进度，RabbitMQ 手工 ACK 则确认某个队列中的投递责任。Kafka 日志因保留期限或容量被清理后，过度落后的组要根据原事实补取或核查缺失输入~\cite{kafka-design,kafka-consumer-api}。

图~\ref{fig:ch10-kafka-groups} 取 G 已完成 100 并提交 101、审计组尚未完成 100 的时刻。两组沿自己的进度读取同一日志，G 内部换成 W2 以后，恢复依据来自该组已保存的位置。

\CTenFigure{ch10-kafka-groups}{传统消费组各自保存分区进度，实例重新分配沿原组位置恢复}{fig:ch10-kafka-groups}

再从 G 的 \code{next=100} 初态看并发处理。W1 取回 100、101、102，并发交给工作线程。100 和 102 已在履约库完成，101 的事务阻塞。若见到 102 完成就为 G 提交 103，重新取得分区的 W2 会从 103 继续。101 位于已提交位置以前，却从未完成，形成遗漏。

安全边界是实际取回记录的连续完成前缀。在这条历史中只能提交 101。后来恢复者可能重读 101 和 102：101 按原身份继续，102 恢复已保存结果，再把位置推进到 103。Kafka 的官方消费者文档也要求并行处理时协调分区内已完成工作的提交边界~\cite{kafka-consumer-api}。

\CTenFigure{ch10-original07}{RabbitMQ 累计确认和 Kafka 分区进度都可能越过尚未完成的工作}{fig:ch10-cumulative}

前缀按当前分配内实际取得的记录与读取位置判断，无需假定每个整数都有应用消息。日志压缩保留原 offset，却可以移除某些记录。事务控制记录和被过滤的中止事务也会造成数值空隙~\cite{kafka-design,kafka-consumer-api}。若实际取回的是 100、102，二者都完成且读取位置允许推进到 103，就可以提交 103。

\CTenFigure{ch10-kafka-prefix}{完成前缀沿实际取得的记录推进，数值空隙与未完成记录分别判断}{fig:ch10-kafka-prefix}

实现可以按分区串行处理，或在一个受控协调器中保存已取回批次、当前分配身份及每项完成状态。关闭自动提交，暂停积压分区、限制并发，只有当前分配有权推进对应位置。提交报错或结局未知时，查询或恢复原进度并容纳重读，履约库按 I/B 保护效果。

分区被重新分配时，旧处理器已经发出的数据库调用或仓储调用可能在途。组协调决定谁继续取得和提交分区进度，另一系统已经开始的动作则由其接收端处理。停止派发新工作、按分配协调未决项，有助于控制恢复。跨系统的效果依据自己的稳定身份和提交规则合流。

\subsection*{Kafka 事务覆盖哪些输出}

前面的消费历史把处理结果写入履约库。现在进一步把输出改成另一 Kafka 主题中的结果记录，输入与输出都由 Kafka 管理，就可以用事务生产者把输出记录与传统组 G 的消费 offset 共同提交。事务生产者配置稳定的 \code{transactional.id}，并发逻辑生产者各自使用不同值，重启时恢复自己的原值。消费者关闭自动提交，输出读取者采用 \code{read\_committed}。offset 提交带当前消费组的关联依据，事务与重平衡按所用 API 的规则协调~\cite{kafka-design,kafka-consumer-api}。

事务提交后，输出与进度共同成立。中止时，相关输出对 \code{read\_committed} 读取者不可见，持久消费位置保持在原提交处。默认 \code{read\_uncommitted} 可以看到中止事务的记录，因此读取隔离也是这项结果的条件~\cite{kafka-design}。

中止以后，当前消费者的内存位置保持着先前读取造成的变化。\footnote{Kafka 4.1 Design 对中止后的当前位置指出 ``does not automatically rewind''。应用恢复时还要处理当前位置和未提交计算状态。}例如原持久位置为 100，poll 后当前读到 103，随后事务中止。若只调用 \code{abortTransaction} 就继续下一次 poll，应用可能跳过本应重新处理的输入。恢复路径丢弃未提交计算状态，按有效的分区分配取得最后已提交位置并显式 seek，或重新建立生产者和消费者、重新分配并从已提交位置读取。旧分配失效时按新分配恢复~\cite{kafka-design}。

同一个 Kafka 事务内调用 PostgreSQL，PostgreSQL 的提交由自己的事务决定。若先提交 F1，再提交 Kafka 输出和 offset，Kafka 中止以后输入会重读，F1 已经存在，处理器按本章 I/B 恢复。若先推进 Kafka 进度，再提交 F，退出又会遗漏履约。Kafka 官方对外部目的系统要求输出与消费进度的额外协作~\cite{kafka-design}。

可以把分区位置与本地输出共同存进履约库，再把该位置作为恢复依据，也可以像本章这样保留稳定业务事件的处理记录，允许代理重读。前一种方案同样需要完成前缀、分区归属及旧在途写入的规则。生产者按 ID/序号处理协议重试的能力，则沿第九章已经区分的范围工作~\cite{kafka-producer}。应用重放 EV1 后，履约消费者通过 I/B 判断已有结果。

将两套输入模型与前面的回报、发货放在图~\ref{fig:ch10-boundaries} 中比较，各个提交域保持自己的恢复依据。Kafka 输入/输出事务覆盖其管理的记录，履约库 Q 和外部发货则在各自边界确认业务效果。

\CTenFigure{ch10-original08}{Kafka、履约库、购买库与外部发货各有自己的提交边界}{fig:ch10-boundaries}
